\chapter{Simple Github}\label{chapter: Simple GitHub}

Simple GitHub usage is when we look around the code repository, or click on a link, and we then simply download the code for our own use. The salient points are:

\begin{itemize}
    \item We don't need (or want) an account on GitHub. We certainly do not wish to login.
    \item We don't intend to amend the code, just use it, although we might want to read through it for our own edukashun\footnote{It's like an education, but better! ;o)}.
    \item If we did change the code in any way, we would not be able to upload it to GitHub to the original authors.
\end{itemize}

\section{What to Do}

We will either have been supplied with a link to the code repository on GitHub, or, we will know what to search for.

\begin{itemize}
    \item Connect to \url{https://github.com} or to the URL supplied---for example, \url{https://github.com/NormanDunbar/QLAssemblyLanguageMagazine}.
    \item Click the green ``<> Code'' button.
    \item Click the ``Download ZIP'' option.
    \item Wait \ldots
    \item Unzip and use the code as desired---bear in mind that you might have to compile the code if the download does not contain a compiled binary. Those normally live in ``Releases'' but can sometimes also be found in the main repo itself.
    \item If the binary you wanted is not in the repo, check for any releases as discussed in \chapref{chapter: Very Simple GitHub}, and download the binary from the releases area.
\end{itemize}